\section{在线策略 vs. 离线策略}

\begin{knowledgebox}{On-policy 与 Off-policy 的区别}
\begin{itemize}
    \item \textbf{离线策略（Off-policy）}：成对比较数据不是来自当前模型的输出。例如 Tulu 3 中使用多个不同模型生成回复，再用 LM 评判。数据一次性收集完毕后反复使用。
    \item \textbf{在线策略（On-policy）}：成对比较数据来自当前正在训练的模型的最新输出。PPO 本质上是 on-policy 方法，需要模型实时生成输出并获取反馈。
\end{itemize}
Off-policy 数据告诉你关于``你不在的地方''的偏好景观；On-policy 数据帮助你``精炼自己''。Tulu 3 同时使用了两种数据。
\end{knowledgebox}

DPO 通常以 off-policy 方式使用（因为数据在训练前就已收集完毕），而 PPO 天然是 on-policy 的。这一区别在实践中有重要影响——on-policy 方法通常能更好地避免 SFT 中讨论的幻觉问题，因为它只让模型在自己的知识范围内优化。


